\chapter{Общий цепно-Hodge след и граница единственности метрики}
\label{ch:v7-common-chain-number-hodge-relative-trace-gate}

\section{От двух положительных норм к одному следу}

Предыдущий гейт вывел относительную производную
\begin{equation}
 \delta_N(F_A)=\frac12[N,F_A],
 \qquad
 \|\delta_N(F_A)\|_{\rm HS}^2=2\|\mathcal R_U(A)\|_{\rm HS}^2.
 \label{eq:v7-common-relative-component}
\end{equation}
Оставался вопрос, можно ли считать рёберный Hodge-момент и эту
relative-компоненту одним физическим следом, а не двумя независимо
нормированными функционалами.

На исследуемом 27-мерном вещественном срезе рёберную кривизну удобно
представить самосопряжённым удвоенным блоком
\begin{equation}
 \mathbb M_E=\operatorname{diag}(m_E,-m_E),
 \qquad \dim_{\mathbb C}\mathbb M_E=54,
 \label{eq:v7-common-edge-curvature}
\end{equation}
где координаты $m_E$ равны
\begin{equation}
 \sqrt{m_a}(|z_a|^2-1),\qquad
 |w_d|^2+\frac85,\qquad
 |w_W|^2+\frac9{10}.
 \label{eq:v7-common-edge-components}
\end{equation}
Тогда
\begin{equation}
 \frac12\Tr_{54}\mathbb M_E^2
 =\mathcal S_E+\text{постоянная}.
 \label{eq:v7-common-edge-trace}
\end{equation}

Связывающий блок сделаем самосопряжённым умножением антиэрмитова
коммутатора на $i$:
\begin{equation}
 \mathbb M_L=i\,\delta_N(Q_A^2-Q_{A_0}^2),
 \qquad \dim_{\mathbb C}\mathbb M_L=42.
 \label{eq:v7-common-linking-curvature}
\end{equation}

\section{Один 96-мерный кривизностный носитель}

Положим
\begin{equation}
 \mathcal K_{\rm common}
 =\mathcal K_E\oplus\mathcal H_{\rm link},
 \qquad
 \dim_{\mathbb C}\mathcal K_{\rm common}=54+42=96,
 \label{eq:v7-common-chain-carrier}
\end{equation}
и
\begin{equation}
 \mathbb F_{\rm common}
 =\operatorname{diag}(\mathbb M_E,\mathbb M_L).
 \label{eq:v7-common-curvature}
\end{equation}
Обычный полный матричный след даёт одно действие
\begin{align}
 \mathcal S_{\rm common}
 &=\frac12\Tr_{96}\mathbb F_{\rm common}^2-\text{const}\nonumber\\
 &=\mathcal S_E+\|\mathcal R_U(A)\|_{\rm HS}^2.
 \label{eq:v7-common-action}
\end{align}
На ста случайных конфигурациях максимальная невязка этого тождества меньше
$10^{-8}$. Никакого относительного коэффициента в
\eqref{eq:v7-common-action} не вставлено.

Полная матричная алгебра $M_{96}(\mathbb C)$ здесь является только
контейнером следа. Она не объявляется координатной алгеброй и не создаёт
новую калибровочную группу. Физическое действие остаётся блок-диагональным
на уже построенных носителях.

\section{Real-завершение}

Real-сопряжение добавляет комплексно-сопряжённую копию общей кривизны:
\begin{equation}
 \mathbb F_{\mathbb R}
 =\mathbb F_{\rm common}\oplus\overline{\mathbb F_{\rm common}}.
 \label{eq:v7-common-real-completion}
\end{equation}
Один физический полуслед удовлетворяет
\begin{equation}
 \frac14\Tr_{192}\mathbb F_{\mathbb R}^2
 =\frac12\Tr_{96}\mathbb F_{\rm common}^2.
 \label{eq:v7-common-real-half-trace}
\end{equation}
Максимальная машинная ошибка меньше $10^{-8}$. Real-удвоение не меняет
отношение двух компонент.

\section{Полный локальный гессиан}

Связывающая relative-компонента начинается с четвёртого порядка в нуле.
Поэтому общий след сохраняет рёберный запуск
\begin{equation}
 (n_-,n_0,n_+)_{0}=(7,0,20),
 \qquad
 \lambda_{\rm heavy}^{\min}=\frac{18}{5}.
 \label{eq:v7-common-origin-signature}
\end{equation}
В целевом вакууме единый след даёт
\begin{equation}
 (n_-,n_0,n_+)_{A_0}=(0,0,27),
 \qquad
 \lambda_{\min}=4.2355904499\ldots.
 \label{eq:v7-common-vacuum-signature}
\end{equation}

Более сильный контроль состоит в замене linking-метрики на любой
положительный вес $\eta\ge0$. Для проверенных значений
\begin{equation}
 \eta=0,\frac14,1,4,100
 \label{eq:v7-common-positive-weight-scan}
\end{equation}
нулевая сигнатура остаётся $(7,0,20)$, а вакуумная --- $(0,0,27)$.
Фактически это следует из двух структурных фактов:
\begin{enumerate}
 \item linking-компонента имеет нулевой гессиан в начале;
 \item её вакуумный гессиан положительно полуопределён.
\end{enumerate}
Следовательно, качественный селектор больше не зависит от относительной
положительной нормировки.

\section{Почему количественная метрика ещё не единственна}

Несмотря на существование одного незавешенного следа, два подносителя имеют
разные размерности:
\begin{equation}
 \dim\mathcal K_E=54,\qquad
 \dim\mathcal H_{\rm link}=42.
 \label{eq:v7-common-dimension-mismatch}
\end{equation}
Между ними нет унитарной обменной симметрии. Проекторы на два ортогональных
слагаемых коммутируют со всей текущей блок-диагональной физической
структурой. Поэтому симметрии не запрещают заменить Hodge-метрику на
\begin{equation}
 \|\mathbb M_E\|^2+\eta\|\mathbb M_L\|^2,
 \qquad \eta>0.
 \label{eq:v7-common-metric-family}
\end{equation}
Выбор $\eta=1$ является минимальным незавешенным общим следом, но пока не
теоремой единственности метрики.

\begin{theorem}[Качественное замыкание общим следом]
Рёберный Hodge-момент и относительная цепная кривизна допускают один
96-мерный самосопряжённый кривизностный носитель и один незавешенный
матричный след. Он даёт правильный селективный запуск $(7,0,20)$ и строго
устойчивый вакуум $(0,0,27)$. Более того, эти два качественных результата
сохраняются при любой положительной относительной Hodge-нормировке.
\end{theorem}

\begin{remark}
Тем самым прежний нормировочный барьер понижен: он больше не угрожает
существованию вакуума и выбору правильных направлений. Но количественные
отношения масс зависят от $\eta$, пока два подносителя не получены как
степени одного неприводимого бикомплекса с единственной Hodge-метрикой.
\end{remark}

\begin{nogocriterion}[Запрет преждевременного массового предсказания]
Нельзя использовать незавешенный выбор $\eta=1$ для численного предсказания
масс, пока единственность общей Hodge-метрики не выведена из total-degree,
Real-структуры или неприводимости общего бикомплекса. Допустимо утверждать
только нормировочно независимые сигнатуры и существование локального
вакуума.
\end{nogocriterion}

\begin{protocol}
\begin{enumerate}
 \item рёберный кривизностный носитель: \textbf{$54$};
 \item связывающий носитель: \textbf{$42$};
 \item общий носитель: \textbf{$96$};
 \item один незавешенный след: \textbf{получен};
 \item reduced action: \textbf{$\mathcal S_E+\|\mathcal R_U\|^2$};
 \item Real-удвоение: \textbf{$192$};
 \item физический полуслед: \textbf{отношение сохраняет};
 \item нулевая сигнатура: \textbf{$(7,0,20)$};
 \item вакуумная сигнатура: \textbf{$(0,0,27)$};
 \item устойчивость при всех положительных весах: \textbf{получена};
 \item качественный нормировочный разрыв: \textbf{закрыт};
 \item количественная Hodge-метрика: \textbf{не единственна};
 \item массовые отношения: \textbf{не выведены};
 \item следующий гейт: \textbf{бикомплекс полной степени}.
\end{enumerate}
\end{protocol}

Машинный сертификат сохранён в
\path{s2t_v7_common_chain_number_hodge_relative_trace_gate_results.json}.

\begin{thebibliography}{9}
\bibitem{v7-common-trace-quillen}
D.~Quillen,
\emph{Superconnections and the Chern Character},
Topology 24 (1985), 89--95.
\bibitem{v7-common-trace-morita}
K.~Kodaka, T.~Teruya,
\emph{The Strong Morita Equivalence for Inclusions of $C^*$-Algebras and
Conditional Expectations for Equivalence Bimodules},
arXiv:1609.08263.
\end{thebibliography}